\section{Research Project Experience}

\setlength{\arrayrulewidth}{.01pt}
\arrayrulecolor[HTML]{C1C1C1}
\setlength{\tabcolsep}{8pt}

\underline{NASA DEVELOP Projects}:

\begin{tabular}{r|p{15.5cm}}
    \textsc{June 2022} \textendash
    &\textbf{Pacific Northwest Health and Air Quality} \\
    \\
    & \textsc{\underline{Title:}} Understanding Threats to the Riverside Shrub Virginia Spiraea: Utilizing Botanical Gardens as a Pathway to Native Plant Conservation\\
    \emph{Present} 
    &\footnotesize{\emph{NOAA National Centers for Environmental Information}}\\
    &\footnotesize{Propose, design, and manage research projects in collaboration with NASA and NOAA science advisors, federal agencies, local government, and non-profit organizations that apply NASA Earth observations to environmental and policy decisions.}
    \\&\\

{\textbf{Pacific Northwest Health and Air Quality}}{NOAA NCEI}
    {\textsc{Role:} Project Manager}{June\textendash Aug. 2023}
    \innerlist{
        \entry{\underline{\textsc{Advisors:}} Dr. Jesse Bell, Dr. Jesse Berman, Dr. Babak Fard, Ronald Leeper, Jared Rennie, Molly Woloszyn}
        \entry{\underline{\textsc{Project Team:}} Greta Bolinger, Tallis Monteiro, Abby Sgan, Cristina Villalobos-Heredia, Taylor West}
        \entryextra{Analyzed trends in PM2.5, PM10, and aerosol optical depth during drought conditions in Oregon and Washington in collaboration with local health departments to support public health response to drought events.}
}

\end{tabular}